\section*{Capítulo \ref{ch:g12:destruction-of-mountains} --- La destrucción de las montañas}

\begin{solution}{exo:g12:destruction-of-mountains:1}
(a) la \omterm{def:g12:destruction-of-mountains:denudation}{tasa de denudación}; (b) la \omterm{def:g12:destruction-of-mountains:load}{carga de fondo}; (c) la \omterm{def:g12:destruction-of-mountains:load}{carga disuelta};
(d) la \omterm{def:g12:destruction-of-mountains:diagenesis}{diagénesis}; (e) el \omterm{def:g12:destruction-of-mountains:rebound}{rebote isostático}.
\end{solution}

\begin{solution}{exo:g12:destruction-of-mountains:2}
(1) Un conglomerado. (2) Una arenisca. (3) Una lutita. (4) Una brecha.
\end{solution}

\begin{solution}{exo:g12:destruction-of-mountains:3}
\qty{1}{m/Myr} son \qty{1000}{mm} en \num{1000000} años, \qty{0.001}{mm} al
año. (1) \qty{0.054}{mm} al año. (2) \qty{0.218}{mm} al año.
(3) $\qty{10}{mm}$ al año son \qty{10}{m} en 1000 años, \qty{10000}{m/Myr}.
\end{solution}

\begin{solution}{exo:g12:destruction-of-mountains:4}
La carga sólida es el \qty{85}{\percent} del total, así que el total es
$60/0.85
\approx 70.6$ millones de toneladas, y la \omterm{def:g12:destruction-of-mountains:load}{carga disuelta}, unos
$70.6 - 60 \approx
10.6$ millones de toneladas al año.
\end{solution}

\begin{solution}{exo:g12:destruction-of-mountains:5}
Meteorización del granito, transporte por un río, sedimentación en una
cuenca, enterramiento bajo capas más jóvenes, cementación.
\end{solution}

\begin{solution}{exo:g12:destruction-of-mountains:6}
A \qty{30}{cm/s}, el grano de arena se levanta (por encima de
\qty{20}{cm/s}); el grano de limo sigue en suspensión. A \qty{5}{cm/s}, el
grano de arena se deposita (por debajo de \qty{10}{cm/s}), pero el grano de
limo sigue en suspensión, porque \qty{0.1}{cm/s} basta para mantenerlo.
\end{solution}

\begin{solution}{exo:g12:destruction-of-mountains:7}
Carga total $4/0.85 \approx 4.71$ millones de toneladas, $\num{4.71e9}/2830 \approx
\qty{1.66e6}{m^3}$ de roca al año. Sobre $\qty{5e10}{m^2}$: $\num{3.3e-5}$ m
al año, unos \qty{33}{m/Myr}, por debajo de la mediana de \qty{54}{m/Myr}.
\end{solution}

\begin{solution}{exo:g12:destruction-of-mountains:8}
(1) $1 - 2700/3300 \approx 0.182$. (2) Unos \qty{182}{m}. (3) Algo más que
los \qty{163}{m} del capítulo: cuanto más ligera es la corteza respecto al
manto, menos sube la columna.
\end{solution}

\begin{solution}{exo:g12:destruction-of-mountains:9}
(1) \qty{8}{km} de roca. (2) $8000 \times 0.163 \approx \qty{1300}{m}$.
(3) $8000/200 = 40$ millones de años.
\end{solution}

\begin{solution}{exo:g12:destruction-of-mountains:10}
(1) $794/43 \approx 18$ millones de años. (2) La superficie solo baja $0.163$
veces el espesor erosionado, así que la roca que hay que quitar es $794/0.163
\approx \qty{4900}{m}$, lo que lleva unos $4900/43 \approx 110$ millones de
años. (3) Los continentes no solo se desgastan: la formación de montañas y el
magmatismo de los arcos no dejan de elevarlos y de añadirles corteza nueva.
\end{solution}

\begin{solution}{exo:g12:destruction-of-mountains:11}
(1) $\num{1e14}/2830 \approx \qty{3.5e10}{m^3}$, es decir, unos
\qty{35}{km^3}. (2) $100/15.5 \approx 6.5$ años.
\end{solution}

\begin{solution}{exo:g12:destruction-of-mountains:12}
$E = kh$ y $h = h_0 e^{-fkt}$ dan $E(t) = kh_0 e^{-fkt}$. El espesor
quitado de $t = 0$ a $t \to \infty$ es $\int_0^{\infty} kh_0
e^{-fkt}\,\mathrm{d}t = kh_0/(fk) = h_0/f$. Para bajar la superficie $h_0$,
la erosión debe quitar $h_0/f \approx 6.1\,h_0$, como dice la proposición del
rebote: cada metro erosionado solo baja la superficie $f$ metros.
\end{solution}

\begin{solution}{exo:g12:destruction-of-mountains:13}
(1) $10 \times 1000 = \qty{10000}{mm}$, \qty{10}{m}. (2) Quitar una carga de
hielo de espesor $t$ equivale a quitar una masa $\rho_{\text{hielo}}\,t$ por
unidad de superficie; la columna sube hasta desplazar esa misma cantidad
menos de manto: $2000 \times 900/3380 \approx \qty{530}{m}$. (3) El terreno
sube cada vez menos a medida que se acerca a su nuevo equilibrio: lo que
impulsa la subida es la parte aún no recuperada.
\end{solution}

\begin{solution}{exo:g12:destruction-of-mountains:14}
(1) Antes de 1900: unos \qty{56}{m/Myr} (como en el ejemplo). En
1987--2006: $145/0.85 \approx 171$ millones de toneladas, $\num{1.71e11}/2830
\approx \qty{6.0e7}{m^3}$, sobre $\qty{2.98e12}{m^2}$: unos
\qty{20}{m/Myr}. (2) La primera: la carga posterior está retenida por presas
y otras obras. (3) No: la disminución es de $400 - 145 = 255$ millones de
toneladas al año, de las que las presas explican de 100 a 150, más o menos la
mitad.
\end{solution}

\begin{solution}{exo:g12:destruction-of-mountains:15}
(1) Al menos \qty{6}{km}. (2) $310 - 252 = 58$ millones de años como máximo,
así que al menos $6000/58 \approx \qty{100}{m/Myr}$. (3) Está entre la
mediana (\qty{54}{m/Myr}) y la media (\qty{218}{m/Myr}) de las cuencas: una
tasa corriente para un cinturón montañoso.
\end{solution}

\begin{solution}{pb:g12:destruction-of-mountains:1}
\textbf{1.} $E = 0.209 \times 2000 \approx \qty{418}{m/Myr}$.

\textbf{2.} \qty{0.418}{mm} al año.

\textbf{3.} Más rápida que ambas: es una cordillera empinada y joven, tan
rápida como las cuencas sísmicamente activas, o más.

\textbf{4.} Su pendiente media.

\textbf{5.} Una capa de $\qty{0.418e-3}{m}$ sobre $\qty{1e6}{m^2}$ son
\qty{418}{m^3}, $418 \times 2830 \approx \qty{1.18e6}{kg}$: unas 1200
toneladas.

\textbf{6.} $2000/418 \approx 4.8$ millones de años.

\textbf{7.} $h(t) = 2000\,e^{-0.209\,t}$ ($t$ en millones de años).

\textbf{8.} $1/0.209 \approx 4.8$ millones de años.

\textbf{9.} $\ln 2/0.209 \approx 3.3$ millones de años.

\textbf{10.} $e^{-0.209 \times 11} = e^{-2.30} \approx 0.100$.

\textbf{11.} $1 - 2830/3380 \approx 0.163$.

\textbf{12.} La superficie baja $100 \times 0.163 \approx \qty{16}{m}$; la
\omterm{def:g12:continental-crust:root}{raíz} sube unos \qty{84}{m}.

\textbf{13.} $h' = -0.163 \times 0.209\,h \approx -0.0341\,h$, así que $h(t) =
2000\,e^{-0.0341\,t}$.

\textbf{14.} $(2000 - 200)/0.163 \approx \qty{11000}{m}$ de roca.

\textbf{15.} Rocas que estaban a unos \qty{11}{km} de profundidad: rocas
metamórficas como esquistos y gneises, y granitos.

\textbf{16.} $1/0.0341 \approx 29$ millones de años.

\textbf{17.} $29 \times \ln 10 \approx 67$ millones de años.

\textbf{18.} $2000\,e^{-0.0341 \times 200} \approx 2000 \times 0.0011 \approx
\qty{2}{m}$.

\textbf{19.} $2000\,e^{-0.0341 \times 250} \approx 2000 \times 0.0002 \approx
\qty{0.4}{m}$: una llanura. Una cumbre de \qty{2037}{m} significa que el
modelo por sí solo no encaja: la región tiene que haberse elevado de nuevo
desde entonces, o haberse erosionado más despacio de lo que supone el modelo.

\textbf{20.} La altitud media se reduce a la mitad en $29 \times \ln 2 \approx$
\textbf{20 millones de años}.
\end{solution}
